% UBT-AI-PROVENANCE-BEGIN
% schema: ubt-ai-provenance/v1
% tier: C_working
% ai_assistance: disclosed
% human_review: risk-based
% editorial_responsibility: Ing. David Jaroš
% policy: ../../../AI_PROVENANCE.md
% notice: Working material; exhaustive human review is not claimed.
% UBT-AI-PROVENANCE-END

\chapter{Kovariantní derivace, konexe a čistá geometrie UBT}
\label{ch:covariant-tetrad}

Tato kapitola vysvětluje geometrické jádro UBT ve formě nejbližší
původní intuice teorie:
\[
 \Theta\quad\longrightarrow\quad E_\mu:=D_\mu\Theta
 \quad\longrightarrow\quad
 \{E_\mu,E_\nu\}\ \hbox{and}\ [D_\mu,D_\nu].
\]
Antikomutátor nese metriku, zatímco komutátor kovariance
deriváty mají zakřivení.  Žádná vkládací mapa, kompaktní vlákno průměr, popř
Projekce pozorovatele je nutná v minimální místní konstrukci.

\section{Obyčejná, parciální a kovariantní derivace}
Pro běžnou skalární funkci $f(x)$ derivace $df/dx$ měří, jak
číslo $f$ se změní, když se změní $x$.  Pro pole na časoprostoru používáme částečné
deriváty,
\[
 \partial_\mu\Theta:=\frac{\partial\Theta}{\partial x^\mu},
 \qquad \mu=0,1,2,3.
\]
Částečná derivace porovnává komponenty zapsané na pevné bázi.  Tohle je
dostačující, lze-li všude použít stejný základ.  Geometrie se stává jemnější
když se samotný lokální snímek otáčí nebo zesiluje z bodu do bodu.

Jednoduchá analogie je vektor zapsaný v rotujícím souřadnicovém rámci.  Dokonce a
fyzikálně konstantní vektor získává měnící se složky, protože základ je
měnící se.  Kovariantní derivát koriguje tuto základní změnu:
\[
 D_\mu\Theta=\partial_\mu\Theta+\rho_*(\Omega_\mu)\Theta.
\]
Zde je $\Omega_\mu$ spojení a $\rho_*$ nám říká, jak spojení
vystupuje v zastoupení $\Theta$.  Pro biquaterniony levá akce, a
správná akce a oboustranná akce jsou skutečně odlišné a nesmí být
tiše vyměněny.

\section{Co je konexe?}
Spojení není silou samo o sobě.  Je to pravidlo pro porovnávání lokálních snímků
v sousedních bodech.  Symbol $\Omega_\mu$ je interní nebo UBT
bikvaternionické spojení rámu.

symboly Christophera,
\[
 \Gamma^\rho{}_{\mu\nu},
\]
spolu úzce souvisejí, ale působí na souřadnicové indexy.  Například,
\[
 \nabla_\mu V^\rho=\partial_\mu V^\rho
 +\Gamma^\rho{}_{\mu\sigma}V^\sigma.
\]
Tedy:
\begin{itemize}
 \item $\Gamma^\rho{}_{\mu\nu}$ nám říká, jak se mění souřadnicová báze;
 \item $\omega_\mu{}^a{}_b$ a jeho spin/bikvaternionový výtah $\Omega_\mu$,
       řekněte nám, jak se mění místní Lorentzův rámec.
\end{itemize}
Jsou spojeny identitou kompatibility tetrád
\[
 \partial_\mu e_\nu{}^a-\Gamma^\rho{}_{\mu\nu}e_\rho{}^a
 +\omega_\mu{}^a{}_b e_\nu{}^b=0.
\]
Jsou to tedy dva popisy stejného geometrického transportu, nikoli
stejné koeficienty.  Zejména $\Gamma=\operatorname{Re}\Omega$ není a
platnou identifikaci.

\section{Proč klasická konexe není další volné pole}
Pro nedegenerovanou tetrádu určují metrická kompatibilita a nulová torze a
unikátní Lorentzovo spojení:
\[
 \boxed{
 \mathring\omega_\mu{}^a{}_b
 =e_b{}^\nu\left(\mathring\Gamma^\rho{}_{\mu\nu}(g)e_\rho{}^a
 -\partial_\mu e_\nu{}^a\right).}
\]
Kruh označuje spojení Levi--Civita bez torze.  Tedy v
obyčejná bezkrutná odbočka GR, spojení rámu se vypočítává z
tetráda, stejně jako se Christoffelovy symboly počítají z metriky.  Místní
Lorentzova rotace mění číselné koeficienty tetrády a tetrády
spojení, ale nemění geometrii.

\subsection{Trať a počítadlo}
Torzní dvoutvar je
\[
 T^a=de^a+\omega^a{}_b\wedge e^b.
\]
Lze zapsat obecné metricky kompatibilní spojení
\[
 \boxed{\omega=\mathring\omega(e)+K(T),}
\]
kde $K$ je zkroucení.  S
$T^a=\tfrac12T^a{}_{bc}e^b\wedge e^c$ a výše uvedená konvence,
\[
 \boxed{
 K_{abc}=\frac12\left(T_{cab}-T_{abc}-T_{bca}\right).}
\]
Tento vzorec má důležitý koncepční důsledek: jednou tetráda a
jsou specifikovány kroucení, metricky kompatibilní spojení je jedinečné.  Otevřené
otázka full-UBT tedy není libovolným zbytkem $\Omega_\mu$; to je
dynamický zákon pro výběr $T$.

Pokud rovnice vakua UBT implikují $T=0$, spojení se redukuje na
Levi--Civita spin připojení.  Pokud se vytvoří spin nebo jiný bikvaternionový zdroj
kroucení, stejná věta rekonstruuje spojení z tohoto kroucení.

\section{Plochý prostor a explicitní UBT reprezentant}
V plochém Minkowského prostoročasu pomocí inerciálních kartézských souřadnic a konstanty
místní rám, lze si vybrat
\[
 \Gamma^\rho{}_{\mu\nu}=0,
 \qquad \Omega_\mu=0,
 \qquad D_\mu=\partial_\mu.
\]
Toto tvrzení lze posílit z výběru měřidla na explicitní
jednopolní reprezentovat.  Snadný
\[
 \boxed{
 \Theta_{\mathrm{SR}}(x)=\Theta_0+\sqrt{\mathcal N_0}
 \left(i x^0\mathbf1+x^k\mathbf e_k\right).}
\]
Pak
\[
 \mathcal N_0^{-1/2}\partial_0\Theta_{\mathrm{SR}}=i\mathbf1,
 \qquad
 \mathcal N_0^{-1/2}\partial_k\Theta_{\mathrm{SR}}=\mathbf e_k.
\]
Centrálním antikomutátorem těchto čtyř konstantních směrů je Minkowski
metrický.  Každá druhá časoprostorová derivace $\Theta_{\mathrm{SR}}$ zmizí,
takže řeší i standardní bezzdrojový konstantní koeficient druhého řádu
hlavní operátor v této ploché větvi.

Obecněji řečeno, pro jakoukoli konstantní Lorentzovu tetrádu $\bar E_\mu$,
\[
 \Theta_{\mathrm{aff}}(x)=\Theta_0+
 \sqrt{\mathcal N_0}\,\bar E_\mu x^\mu
\]
generuje přesně tu tetrádu.  Tím se uzavírá speciální relativismus
reprezentativní problém, ačkoli to nedokazuje jedinečnost nebo zakřivený prostor
existence.

V polárních souřadnicích nebo v urychlujícím snímku mohou $\Gamma$ a $\Omega$
být nenulový, i když je časoprostor plochý.  Invariantní prohlášení o plochosti
je mizivé zakřivení, nikoli mizející koeficienty spojení v každém snímku.

\section{Čtyři kovariantní gradienty jako tetráda}
Definovat
\[
 E_\mu:=\frac{1}{\sqrt{\mathcal N_0}}D_\mu\Theta.
\]
Čtyři objekty $E_0,E_1,E_2,E_3$ jsou místní pravítka UBT a směr hodin.
Nejsou to čtyři extra hmotná pole; jsou to první kovariantní deriváty
jediného základního pole $\Theta$.

Nechť $i$ je běžná komplexní jednotka pro dojíždění a
$\mathbf e_1,\mathbf e_2,\mathbf e_3$ kvaternionové jednotky.  V klasickém
Lorentz sektor napsat
\[
 E_\mu=i e_\mu{}^0\mathbf1+e_\mu{}^k\mathbf e_k,
 \qquad e_\mu{}^a\in\mathbb R.
\]
Konjugace kvaternionů mění znaky tří kvaternionových jednotek, ale
ponechává komplexní skalární koeficient nedotčený:
\[
 E_\mu^\sharp=i e_\mu{}^0\mathbf1-e_\mu{}^k\mathbf e_k.
\]

\section{Proč je antikomutátor metrikou}
Kvaternionové násobení obsahuje skalární součin a křížový součin.  V
symetrický součin, který křížové produkty ruší:
\[
 \boxed{
 \frac12\left(E_\mu^\sharp E_\nu+E_\nu^\sharp E_\mu\right)
 =g_{\mu\nu}\mathbf1.}
\]
Pravá strana obsahuje $\mathbf1$, jednotku biquaternion.  celý
levá strana je již centrální: obyčejné reálné číslo vynásobené
jednotka algebry.  Žádná stopa, mapa skutečných částí, fázová projekce nebo průměr vlákna
požadovaný.

Rozšíření koeficientů dává
\[
 g_{\mu\nu}
 =-e_\mu{}^0e_\nu{}^0+
  e_\mu{}^1e_\nu{}^1+
  e_\mu{}^2e_\nu{}^2+
  e_\mu{}^3e_\nu{}^3,
\]
což je přesně
\[
 g_{\mu\nu}=e_\mu{}^a e_\nu{}^b\eta_{ab},
 \qquad \eta_{ab}=\operatorname{diag}(-1,1,1,1).
\]
Minusové znaménko času pochází z $i^2=-1$.

\section{Co dělá druhá polovina součinu?}
Antisymetrická algebraická část je
\[
 \Sigma_{\mu\nu}
 =\frac12\left(E_\mu^\sharp E_\nu-E_\nu^\sharp E_\mu\right).
\]
Popisuje orientovanou rovinu nebo bivektor a je přirozeným nosičem pro místní
informace o rotaci, zesílení a rotaci.  To samo o sobě není zakřivení.

Tři výrazy musí být odděleny:
\begin{align*}
 \{E_\mu,E_\nu\}_\sharp
 &\longrightarrow \text{metric},\\
 [E_\mu,E_\nu]_\sharp
 &\longrightarrow \text{algebraic bivector/spin structure},\\
 [D_\mu,D_\nu]
 &\longrightarrow \text{connection curvature or gauge field strength}.
\end{align*}

\section{Křivost z nekomutujících derivací}
Pro jednostranné spojení má zakřivení známý tvar
\[
 \mathcal R_{\mu\nu}
 =\partial_\mu\Omega_\nu-\partial_\nu\Omega_\mu
  +[\Omega_\mu,\Omega_\nu].
\]
Koncový komutátor je nezbytný, protože biquaternionové násobení je
nekomutativní.  Geometricky, zakřivení měří, jak se rám poté změní
přepravovat kolem malé uzavřené smyčky.

Rovnice $E_\mu=D_\mu\Theta/\sqrt{\mathcal N_0}$ vytváří další
stav integrovatelnosti.  Toto je bod, ve kterém je reprezentace
spojení se stává rozhodujícím.

\section{Proč starý problém s hodností lokálně mizí}
Symetrická metrika $4\times4$ má deset nezávislých složek.  Při pohledu pouze na
jedna hodnota $\Theta\in\mathbb B\simeq\mathbb R^8$ naznačuje zavádějící
srovnání $8<10$.

Metrika není postavena pouze na hodnotě $\Theta$.  Staví se ze čtyř
kovariantní deriváty $E_\mu$.  V sektoru Lorentz tyto tvoří skutečný
Tetrada $4\times4$$e_\mu{}^a$ se šestnácti součástkami.  Šest změn je lokálních
rotace a zesílení, které ponechávají metriku beze změny.  Proto
\[
 16-6=10.
\]
Přesněji řečeno, každá malá symetrická metrická variace $h_{\mu\nu}$ je
produkoval
\[
 \delta e_\mu{}^a=\frac12 h_{\mu\rho}e^{\rho a}.
\]
Mapa tetrad-to-metric má proto pořadí deset u každé nedegenerované tetrády.
Tím je vyřešena otázka místního kinematického pořadí.  Zatím to neprokazuje
každá zakřivená tetráda je generována polem UBT na skořápce.

\subsection{Správný výpočet může odpovědět na špatně položenou architektonickou otázku}
Dřívější jednodílná hodnostní obstrukce byla matematicky správná v an
indukovaný metrický rámec, kde měly být obnoveny pouze Einsteinovy rovnice
z normálních variant jedné sekce podobné vkládání.  Nepředstavovalo a
no-go teorém pro kovariantní tetrádu $E_\mu=D_\mu\Theta/\sqrt{\mathcal N_0}$.
Konstrukce z kompaktních vláken opravila dřívější složení přidáním
směry profilu, ale formulace tetrád ukazuje, že tato oprava nebyla
vyžaduje původní architektura.

Toto rozlišení motivuje obecné pravidlo výzkumu:

\begin{quote}
Před přidáním polí, režimů, kót, projekcí, průměrů nebo pomocných prvků
struktury k vyřešení překážky, nejprve otestujte, zda se jedná o překážku
artefakt formulace, ze které byl odvozen.
\end{quote}

Trasa vlákna zůstává matematicky konzistentní ve svém vlastním rámci.  Jeho
problém je slabý výběr a redundance zástupce, nikoli algebraika
rozpor.

\section{Problém integrability}
Jakmile je klasické spojení rekonstruováno z tetrády a torze,
definující rovnice se stává schematickou
\[
 E_\mu=\mathcal N_0^{-1/2}
 \left(\partial_\mu\Theta+\rho_*(\Omega_\mu[E,T])\Theta\right).
\]
Neznámý $E$ se vyskytuje na obou stranách.  Toto je implicitní nelineární první řád
PDE nebo systém pevných bodů.  Kinematická eliminace $\Omega$ neprobíhá
sám dokazuje existenci nebo jedinečnost páru $(\Theta,E)$.

\subsection{Jednostranná překážka vedoucí k plochosti}
Předpokládejme, že spojení funguje pouze zleva,
\[
 D^L_\mu\Theta=\partial_\mu\Theta+A_\mu\Theta,
\]
aby
\[
 [D^L_\mu,D^L_\nu]\Theta=F^L_{\mu\nu}\Theta.
\]
Ve větvi bez torzní kompatibilní tetrády, antisymetrizující kovariant
Hessian dává
\[
 F^L_{\mu\nu}\Theta=0.
\]
Pokud je $\Theta$ invertibilní jako biquaternion, pak
\[
 F^L_{\mu\nu}=0.
\]
Naivní jednostrannou regulérní reprezentaci tedy nelze obecně podporovat
zakřivený GR při zachování invertibilního $\Theta$ a torzní tetrády
kompatibilita.  Konfigurace neinvertovatelného stabilizátoru se mohou tomuto výsledku vyhnout,
nelze však předpokládat, že představují libovolné geometrie.

\subsection{Přirozená oboustranná bikvaternionová derivace}
Biquaterniony přirozeně připouštějí levé i pravé násobení.  Zvážit
\[
 \boxed{
 D_\mu\Theta=\partial_\mu\Theta+A_\mu\Theta-\Theta B_\mu.}
\]
Přímé rozšíření dává
\[
 \boxed{
 [D_\mu,D_\nu]\Theta
 =F^A_{\mu\nu}\Theta-\Theta F^B_{\mu\nu}.}
\]
Proto je nutná kompatibilita bez kroucení
\[
 F^A_{\mu\nu}\Theta=\Theta F^B_{\mu\nu}.
\]
Pro invertibilní $\Theta$,
\[
 \boxed{
 F^A_{\mu\nu}=\Theta F^B_{\mu\nu}\Theta^{-1}.}
\]
Na rozdíl od jednostranného pouzdra nesmí zmizet ani zakřivení.  Pole $\Theta$
proplétá levou a pravou reprezentaci zakřivení.  Tohle je opravdový
strukturální zúžení GAP-10I: oboustranný/bimodulový derivát je
minimální algebraně-nativní trasa, která se vyhýbá obecné překážce plochosti.

\subsection{Redukce bez nového pole, bez-torzní no-go a lokální pravá inverze s torzí}
Na Lorentzově plátku $W_L=\{X\mid X^\ddagger=-X\}$, konzervace plátku
a jeho kvadratické formy redukuje čistý pár, modulo společný centrální
jeden formulář, který ruší, do
\[
 \boxed{A_\mu=\Omega_\mu,\qquad B_\mu=-\Omega_\mu^\ddagger.}
\]
$A_\mu$ a $B_\mu$ tedy v tomto nejsou nezávislá gravitační pole
větev.  Jsou to dvě modulové akce připojení s jedním spinem.

S nulovou torzí vytváří stejná involučně kompatibilní derivace a
nedegenerované řešení $D_\mu\Theta=\sqrt{\mathcal N_0}E_\mu$ implikovat
\[
 \boxed{\mathring\nabla_\mu V^\nu=\delta_\mu{}^\nu,}
 \qquad \mathcal L_Vg=2g.
\]
Neplochý Schwarzschildův vakuový exteriér s nenulovou hmotností nic takového nemá
stejnorodost.  Čistý pár je tedy uzavřen kinematicky a jeho $K=0$
větev je uzavřena jako beztorzní no-go.

Na kvalifikaci záleží.  Na dostatečně malém nenulovém Gaussově políčku
zvolte $\rho\ne0$ s
$g^{\mu\nu}\partial_\mu\rho\partial_\nu\rho=\varepsilon=\pm1$ a sada
\[
 V^\mu=\varepsilon\rho\nabla^\mu\rho,
 \qquad
 W_{\mu\nu}=g_{\mu\nu}-\mathring\nabla_\mu V_\nu,
\]
\[
 \boxed{K_{\nu\mu\rho}
 =\frac{W_{\mu\nu}V_\rho-V_\nu W_{\mu\rho}}{V^2}.}
\]
Poté $K_{\nu\mu\rho}=-K_{\rho\mu\nu}$ a
$\nabla^{(\mathring\Gamma+K)}_\mu V^\nu=\delta_\mu{}^\nu$.  Proto
\[
 \Theta=\sqrt{\mathcal N_0}V^a\mathbf u_a
\]
poslouchá $D_\mu\Theta=\sqrt{\mathcal N_0}E_\mu$.  Každá hladká tetráda má takový
lokální zástupce s kompozitním metrickým kompatibilním zkreslením a ne
nezávislá pole $A_\mu,B_\mu$.

Relativní výraz $P_\mu X-XQ_\mu$ již není pro lokální kinematiku vyžadován
existence.  Zůstává možné bez torzní kompletace a je-li použito, musí být
kompozitní nebo pomocné spíše než nové pole šíření.  Společná centrála
konstantní neodstraní překážku bez torze, protože se ruší.

\section{Podmíněná dynamická dílčí uzavření}
\subsection{Algebraická Cartanova rovnice torze}
V minimální větvi Hilbert-Palatini prvního řádu, nezávislá variace
spojení Lorentz dává
\[
 \epsilon_{abcd}T^a\wedge e^b=\kappa\tau_{cd}.
\]
Pro nedegenerovanou tetrádu je to $24\times24$ bodová lineární mapa hodnosti
24. Tedy nulový spinový proud dává $T=0$, zatímco specifikovaný spinový proud dává a
jedinečná torze a zkroucení. To uzavře torzní algebru uvnitř
minimální větev, nikoli odvození této větve od fundamentálního UBT.

\subsection{Propagace Lorentzova řezu jako množiny pevných bodů}
Antilineární involuce
\[
 \mathcal JX=-\overline{X^\sharp}
\]
má fyzický Lorentz modul jako pevnou sadu. Pokud je dobře nasazený unikát
Cauchy evolution je $\mathcal J$-ekvivariantní a jeho data a zdroje jsou pevné
od $\mathcal J$, jedinečnost znamená, že řešení zůstává v Lorentz
plátek. Konečná akce UBT musí být ještě porovnána s těmito hypotézami.

\subsection{Stabilita metriky v imaginárním čase}
Vertikální tok
\[
 \partial_\psi e_\mu{}^a=\Lambda_\psi{}^a{}_b e_\mu{}^b,
 \qquad \Lambda_{\psi ab}=-\Lambda_{\psi ba},
\]
je tečnou k místní orbitě Lorentzova měřidla a vyhovuje
$\partial_\psi g_{\mu\nu}=0$ identicky. Druhým dostatečným mechanismem je a
unikátní $\psi$-překlad-invariantní evoluce s $\psi$-nezávislá iniciála
data. Výběr fyzického mechanismu a vyloučení nestabilních neměřicích režimů
zůstává dynamickým úkolem.

\subsection{Rozšířená holonomie pro předepsané zakřivené koeficienty}
Pro předepsaný $E_\mu,A_\mu,B_\mu$, nehomogenní systém
\[
 \partial_\mu\Theta+A_\mu\Theta-\Theta B_\mu
 =\sqrt{\mathcal N_0}E_\mu
\]
se stává homogenním paralelním transportem pro $Y=(\Theta,1)^T$ v rozšířeném
svazek. Jednohodnotové řešení s počáteční hodnotou $Y_0$ existuje přesně kdy
rozšířená holonomie opravuje $Y_0$; ploché rozšířené zakřivení stačí na
jednoduše připojená doména. Tím se uzavře integrovatelnost předepsaného koeficientu,
nekonzistentní výběr koeficientů UBT.

\subsection{Podmíněný Einsteinův--Lambda koncový bod}
Minimální Palatiniho větev s nulovými výtěžky spinového proudu
\[
 G_{\mu\nu}+\Lambda g_{\mu\nu}=\kappa T_{\mu\nu}.
\]
Podle běžných čtyřrozměrných Lovelockových předpokladů je každý místní přirozený
metrická rovnice druhého řádu bez divergence bez extra lehké geometrie
pole má gravitační část $aG_{\mu\nu}+bg_{\mu\nu}$. Proto
podmíněný nízkoenergetický koncový bod je jedinečný. UBT je musí stále odvodit
předpoklady, koeficienty a hmotné členy z jeho kanonické akce.

\section{Implicitní versus transcendentní rovnice}
Slova\emph{implicitní} a\emph{transcendentální} popisují různé
vlastnosti.

Rovnice je implicitní, když se neznámá vyskytuje uvnitř předpokládané mapy
určit to:
\[
 E=\mathcal F(E,\Theta).
\]
Rovnice je transcendentální, když obsahuje nealgebraické funkce jako např
exponenciální, logaritmická, trigonometrická funkce nebo funkce Jacobi theta:
\[
 E=e^{-E}.
\]
Tetradová rovnice UBT je již implicitní a diferenciální.  Pokud je povoleno
pole $\Theta(q,\tau)$ je omezeno na Jacobi-theta nebo jinou transcendentální hodnotu
třídy funkcí, pak může být celý systém jak implicitní, tak i
transcendentální.  Tyto vlastnosti by měly být uvedeny ve formálním dokumentu
odděleně, i když intuitivně oba ``zapletou neznámo
sám.''

\section{Derivace versus variace}
Derivát jako $D_\mu\Theta$ popisuje, jak se změní konfigurace jednoho pole
přes časoprostor.  Varianta porovnává možné konfigurace v okolí:
\[
 \Theta_\varepsilon=\Theta+\varepsilon\,\delta\Theta.
\]
Variace je
\[
 \delta\Theta=
 \left.\frac{d\Theta_\varepsilon}{d\varepsilon}\right|_{0}.
\]
Při odvozování rovnic pole z akce jsou zapotřebí variace.  jsou
nikoli náhrady za časoprostorové deriváty.  Když připojení závisí na
pole,
\[
 \delta(D_\mu\Theta)
 =D_\mu(\delta\Theta)+\rho_*(\delta\Omega_\mu)\Theta
\]
a musí být zahrnuta variace indukovaného připojení.

\section{Co je uzavřeno a co zůstává otevřené}
Následující výsledky jsou uzavřeny na uvedených úrovních:
\begin{enumerate}
 \item centrální antikomutátorová metrika a rank-ten tetrad mapa;
 \item GAP-10$\Omega$-KIN/GR: rekonstrukce kompatibilního připojení a
       její Levi--Civita bez torzní větve;
 \item GAP-10T-PALATINI (podmíněné): minimální Cartanova torzní rovnice je
       algebraické a invertibilní;
 \item GAP-10L-CONN a GAP-10L-SYM (podmíněné): kompatibilní přepravní a
       jedinečná ekvivariantní evoluce zachovat Lorentzovu výseč;
 \item GAP-10I-SR a GAP-10I-PRESCRIBED: afinní plochý zástupce a
       přesné kritérium rozšířené holonomie pro předepsané zakřivené koeficienty;
 \item GAP-10I-1S: naivní jednostranná invertibilní zakřivená trasa je podmínkou
       ne-jít;
 \item GAP-10I-PAIR-KIN: čistý kompatibilní pár se redukuje na jedno otočení
       spojení;
 \item GAP-10I-PAIR-GR: čistý pár $K=0$ je nepřetržitý souběžný vektor
       pro neplochý Schwarzschildův vakuový exteriér s nenulovou hmotností;
 \item GAP-10I-TORSION-LOCAL: každá hladká tetráda má lokální jedno-$\Theta$
       zástupce s explicitním složeným metrickým kompatibilním zkreslením;
 \item GAP-10D-PALATINI/UNIQUENESS (podmíněné): minimální první řád
       větev a Lovelockovy hypotézy vybrat Einstein--$\Lambda$;
 \item GAP-10$\psi$-KIN/SYM: Lorentzova průtoková nebo translační symetrie
       chrání klasickou metriku podle uvedených hypotéz.
\end{enumerate}

Nerozdělené mezery plné teorie jsou spíše zúženy než plně uzavřeny:
\begin{enumerate}
 \item GAP-10T-SPIN je podmíněně uzavřen pro přímé pevné pozadí
       Palatiniho variace, zatímco GAP-10T-FLAT-NOGO a
       GAP-10T-PAIRING-NOGO jsou uzavřeny jako no-go výsledky;
 \item GAP-10T-DYN: vypočítejte plnou složenou variaci $\Theta$ a odvoďte
       vybranou větev I. řádu, kanonické torzní zrušení popř
       translační/relativní dokončení a normalizace z kanonické UBT;
 \item GAP-10I-CURVED: místní kinematika je uzavřena; odvodit kanonický
       výběr na úrovni akce a fyzická přípustnost kompozitu
       torzní větev nebo volitelná nekroucená relativní bimodulová větev,
       a zavést pravidelnost a globální pokračování
       $(\Theta,E,A,B,T)$;
 \item GAP-10L-DYN: dokažte rovnocennost a dobře prezentovanou jedinečnost pro finále
       úplná dynamika;
 \item GAP-10D: odvodit nízkoenergetické předpoklady Palatini/Lovelock, spojky
       a působení hmoty ze základní teorie;
 \item GAP-10$\psi$: vyberte mechanismus stability a vylučte fyzické
       nestabilní neměřicí režimy.
\end{enumerate}
GAP-B-MASTER a GAP-U2$\Theta$ zůstávají otevřené. Dřívější kompaktní vlákno $\psi$
uzavření zůstává matematicky konzistentním historickým kandidátským dokončením,
ale jeho slabý výběr a redundance zástupců jej udržují mimo minimum
kanonická cesta.
